\section{How the Machine Is
Organised}\label{how-the-machine-is-organised}

Helicase LXR stores its work in a hierarchy. Understanding it makes
Load/Save, Morph and live performance much easier to follow.

\figplaceholder{hierarchy_diagram.png}{The Bank, Scene, Kit and Instrument hierarchy.}{Tree diagram. Bank at the top, containing up to 16 Scenes. Each Scene contains: a Pattern (7 tracks of 128 steps), a Kit made of 6 Instruments (one per voice slot), one Effect, mix and MIDI settings, and chain settings. Show Scene 17 as a dashed box labelled background Scene. Use orange for things stored on the SD card as libraries and grey for things that live inside a Scene.}

\subsection{Bank}\label{bank}

A \textbf{Bank} is a complete set of up to \textbf{16 Scenes}. One Bank
is resident (loaded in memory) at a time. Loading a Bank replaces the
Scenes; saving a Bank writes every Scene back as a folder on the SD
card. Banks can be changed from the Load page or over MIDI (CC0 on the
global channel).

\subsection{Scene}\label{scene}

A \textbf{Scene} is everything needed to play one section of a song.
Each of the 16 Scenes holds:

\begin{itemize}
\tightlist
\item
  a \textbf{Pattern}: 7 tracks of 128 steps, with all step data and step
  automation;
\item
  a \textbf{Kit}: six Instruments, one per voice slot;
\item
  one \textbf{Effect} (off, StFilter or CrumpBit) with its parameters
  and its 16-step effect sequence;
\item
  \textbf{mix settings} per voice (output route, effect send, fader
  mode, Morph endpoints);
\item
  the \textbf{bus compressor} settings;
\item
  \textbf{per-track timing} (length, step scale, shuffle, run mode) and
  per-track \textbf{MIDI channel and note};
\item
  the Scene's \textbf{chain} settings (what plays next).
\end{itemize}

Switching Scene in Performance mode changes all of these together.
Because the Scene owns the sounds as well as the notes, changing Scene
can change the whole sound of the machine.

\subsubsection{Active Scene and played
Scene}\label{active-scene-and-played-scene}

The \textbf{active Scene} is the Scene you are looking at and editing:
it decides which values the Voice and Effect pages show and which Scene
your edits are written to. The \textbf{played Scene} of a track is the
Scene whose Pattern that track is currently playing. Normally they are
the same, but with \emph{per-track Scene playback} (see the Operation
Modes chapter) each track can play a different Scene.

\subsubsection{Scene 17}\label{scene-17}

While a Bank is loading in the background, the music keeps playing from
a hidden copy called \textbf{Scene 17}. It is released when you change
Scene.

\subsection{Kit and Instrument}\label{kit-and-instrument}

An \textbf{Instrument} is the patch of one voice: oscillators, filter,
envelopes, LFO and so on. There are several Instrument \textbf{types}
(Drum, Snare, Cymbal, HiHat and Alkaline), each with its own parameter
set. A \textbf{Kit} is a set of six Instruments. Instruments and Kits
can be saved and loaded individually.

Slots have rules:

\begin{itemize}
\tightlist
\item
  \textbf{Basic} types may be used in any number of slots.
\item
  \textbf{Advanced} types are limited to two per Kit. No current type is
  Advanced.
\item
  The \textbf{HiHat} is a \emph{choke} type: in slot 6 it serves both
  track 6 (closed) and track 7 (open), and the closed hat chokes the
  open hat.
\end{itemize}

\subsection{Effect}\label{effect}

Each Scene owns exactly one \textbf{Effect} on a shared effects bus.
Voices feed it through their own send levels. Loading a Kit or an
Instrument never changes the Effect; loading a Scene or a Bank does.

\subsection{Pattern, track and step}\label{pattern-track-and-step}

A \textbf{Pattern} belongs to a Scene. It has \textbf{7 tracks} of
\textbf{128 steps}. The 128 steps of a track are shown as eight bars of
16 steps on the sequencer buttons. A step carries a trigger and optional
\emph{specials}: note, velocity, condition, roll, timing, rhythm and
fill, plus step \textbf{automation} entries that set parameters on that
step.

\subsection{Morph}\label{morph}

Most Instrument, Effect and Scene parameters store \textbf{two values}:
a Normal value and a Morph value. A \textbf{Morph amount} (0 to 255)
blends between them. There is a Morph amount per voice, one for the
Effect, and a \textbf{global Morph} (\texttt{mrp}) that sets all of them
at once. Morph is the main tool for evolving sounds live. See
\emph{Morph} in the Operation Modes chapter.

\subsection{Edit mask}\label{edit-mask}

Edits fan out to every Scene in the \textbf{voice edit mask}, so several
Scenes with the same structure can be edited together. Pattern data is
never fanned out.

\subsection{Automatic saving}\label{automatic-saving}

Edits are written to the SD card automatically in the background
(\textbf{AutoSave}). You do not have to press Save to keep your work in
the current Bank; Save is for creating named library files. See the
Storage chapter.

\draftnote{The terms ``Kit'' and ``Instrument'' follow the firmware. The names of the Load/Save
types and the exact way the voice edit mask is chosen from the panel need confirming on the
device.}
